% Options for packages loaded elsewhere
% Options for packages loaded elsewhere
\PassOptionsToPackage{unicode}{hyperref}
\PassOptionsToPackage{hyphens}{url}
\PassOptionsToPackage{dvipsnames,svgnames,x11names}{xcolor}
%
\documentclass[
  letterpaper,
  DIV=11,
  numbers=noendperiod]{scrartcl}
\usepackage{xcolor}
\usepackage[margin=1in]{geometry}
\usepackage{amsmath,amssymb}
\setcounter{secnumdepth}{5}
\usepackage{iftex}
\ifPDFTeX
  \usepackage[T1]{fontenc}
  \usepackage[utf8]{inputenc}
  \usepackage{textcomp} % provide euro and other symbols
\else % if luatex or xetex
  \usepackage{unicode-math} % this also loads fontspec
  \defaultfontfeatures{Scale=MatchLowercase}
  \defaultfontfeatures[\rmfamily]{Ligatures=TeX,Scale=1}
\fi
\usepackage{lmodern}
\ifPDFTeX\else
  % xetex/luatex font selection
  \setmonofont[]{DejaVu Sans Mono}
\fi
% Use upquote if available, for straight quotes in verbatim environments
\IfFileExists{upquote.sty}{\usepackage{upquote}}{}
\IfFileExists{microtype.sty}{% use microtype if available
  \usepackage[]{microtype}
  \UseMicrotypeSet[protrusion]{basicmath} % disable protrusion for tt fonts
}{}
\makeatletter
\@ifundefined{KOMAClassName}{% if non-KOMA class
  \IfFileExists{parskip.sty}{%
    \usepackage{parskip}
  }{% else
    \setlength{\parindent}{0pt}
    \setlength{\parskip}{6pt plus 2pt minus 1pt}}
}{% if KOMA class
  \KOMAoptions{parskip=half}}
\makeatother
% Make \paragraph and \subparagraph free-standing
\makeatletter
\ifx\paragraph\undefined\else
  \let\oldparagraph\paragraph
  \renewcommand{\paragraph}{
    \@ifstar
      \xxxParagraphStar
      \xxxParagraphNoStar
  }
  \newcommand{\xxxParagraphStar}[1]{\oldparagraph*{#1}\mbox{}}
  \newcommand{\xxxParagraphNoStar}[1]{\oldparagraph{#1}\mbox{}}
\fi
\ifx\subparagraph\undefined\else
  \let\oldsubparagraph\subparagraph
  \renewcommand{\subparagraph}{
    \@ifstar
      \xxxSubParagraphStar
      \xxxSubParagraphNoStar
  }
  \newcommand{\xxxSubParagraphStar}[1]{\oldsubparagraph*{#1}\mbox{}}
  \newcommand{\xxxSubParagraphNoStar}[1]{\oldsubparagraph{#1}\mbox{}}
\fi
\makeatother

\usepackage{color}
\usepackage{fancyvrb}
\newcommand{\VerbBar}{|}
\newcommand{\VERB}{\Verb[commandchars=\\\{\}]}
\DefineVerbatimEnvironment{Highlighting}{Verbatim}{commandchars=\\\{\}}
% Add ',fontsize=\small' for more characters per line
\usepackage{framed}
\definecolor{shadecolor}{RGB}{241,243,245}
\newenvironment{Shaded}{\begin{snugshade}}{\end{snugshade}}
\newcommand{\AlertTok}[1]{\textcolor[rgb]{0.68,0.00,0.00}{#1}}
\newcommand{\AnnotationTok}[1]{\textcolor[rgb]{0.37,0.37,0.37}{#1}}
\newcommand{\AttributeTok}[1]{\textcolor[rgb]{0.40,0.45,0.13}{#1}}
\newcommand{\BaseNTok}[1]{\textcolor[rgb]{0.68,0.00,0.00}{#1}}
\newcommand{\BuiltInTok}[1]{\textcolor[rgb]{0.00,0.23,0.31}{#1}}
\newcommand{\CharTok}[1]{\textcolor[rgb]{0.13,0.47,0.30}{#1}}
\newcommand{\CommentTok}[1]{\textcolor[rgb]{0.37,0.37,0.37}{#1}}
\newcommand{\CommentVarTok}[1]{\textcolor[rgb]{0.37,0.37,0.37}{\textit{#1}}}
\newcommand{\ConstantTok}[1]{\textcolor[rgb]{0.56,0.35,0.01}{#1}}
\newcommand{\ControlFlowTok}[1]{\textcolor[rgb]{0.00,0.23,0.31}{\textbf{#1}}}
\newcommand{\DataTypeTok}[1]{\textcolor[rgb]{0.68,0.00,0.00}{#1}}
\newcommand{\DecValTok}[1]{\textcolor[rgb]{0.68,0.00,0.00}{#1}}
\newcommand{\DocumentationTok}[1]{\textcolor[rgb]{0.37,0.37,0.37}{\textit{#1}}}
\newcommand{\ErrorTok}[1]{\textcolor[rgb]{0.68,0.00,0.00}{#1}}
\newcommand{\ExtensionTok}[1]{\textcolor[rgb]{0.00,0.23,0.31}{#1}}
\newcommand{\FloatTok}[1]{\textcolor[rgb]{0.68,0.00,0.00}{#1}}
\newcommand{\FunctionTok}[1]{\textcolor[rgb]{0.28,0.35,0.67}{#1}}
\newcommand{\ImportTok}[1]{\textcolor[rgb]{0.00,0.46,0.62}{#1}}
\newcommand{\InformationTok}[1]{\textcolor[rgb]{0.37,0.37,0.37}{#1}}
\newcommand{\KeywordTok}[1]{\textcolor[rgb]{0.00,0.23,0.31}{\textbf{#1}}}
\newcommand{\NormalTok}[1]{\textcolor[rgb]{0.00,0.23,0.31}{#1}}
\newcommand{\OperatorTok}[1]{\textcolor[rgb]{0.37,0.37,0.37}{#1}}
\newcommand{\OtherTok}[1]{\textcolor[rgb]{0.00,0.23,0.31}{#1}}
\newcommand{\PreprocessorTok}[1]{\textcolor[rgb]{0.68,0.00,0.00}{#1}}
\newcommand{\RegionMarkerTok}[1]{\textcolor[rgb]{0.00,0.23,0.31}{#1}}
\newcommand{\SpecialCharTok}[1]{\textcolor[rgb]{0.37,0.37,0.37}{#1}}
\newcommand{\SpecialStringTok}[1]{\textcolor[rgb]{0.13,0.47,0.30}{#1}}
\newcommand{\StringTok}[1]{\textcolor[rgb]{0.13,0.47,0.30}{#1}}
\newcommand{\VariableTok}[1]{\textcolor[rgb]{0.07,0.07,0.07}{#1}}
\newcommand{\VerbatimStringTok}[1]{\textcolor[rgb]{0.13,0.47,0.30}{#1}}
\newcommand{\WarningTok}[1]{\textcolor[rgb]{0.37,0.37,0.37}{\textit{#1}}}

\usepackage{longtable,booktabs,array}
\newcounter{none} % for unnumbered tables
\usepackage{calc} % for calculating minipage widths
% Correct order of tables after \paragraph or \subparagraph
\usepackage{etoolbox}
\makeatletter
\patchcmd\longtable{\par}{\if@noskipsec\mbox{}\fi\par}{}{}
\makeatother
% Allow footnotes in longtable head/foot
\IfFileExists{footnotehyper.sty}{\usepackage{footnotehyper}}{\usepackage{footnote}}
\makesavenoteenv{longtable}
\usepackage{graphicx}
\makeatletter
\newsavebox\pandoc@box
\newcommand*\pandocbounded[1]{% scales image to fit in text height/width
  \sbox\pandoc@box{#1}%
  \Gscale@div\@tempa{\textheight}{\dimexpr\ht\pandoc@box+\dp\pandoc@box\relax}%
  \Gscale@div\@tempb{\linewidth}{\wd\pandoc@box}%
  \ifdim\@tempb\p@<\@tempa\p@\let\@tempa\@tempb\fi% select the smaller of both
  \ifdim\@tempa\p@<\p@\scalebox{\@tempa}{\usebox\pandoc@box}%
  \else\usebox{\pandoc@box}%
  \fi%
}
% Set default figure placement to htbp
\def\fps@figure{htbp}
\makeatother





\setlength{\emergencystretch}{3em} % prevent overfull lines

\providecommand{\tightlist}{%
  \setlength{\itemsep}{0pt}\setlength{\parskip}{0pt}}



 


\KOMAoption{captions}{tableheading}
\makeatletter
\@ifpackageloaded{caption}{}{\usepackage{caption}}
\AtBeginDocument{%
\ifdefined\contentsname
  \renewcommand*\contentsname{Table of contents}
\else
  \newcommand\contentsname{Table of contents}
\fi
\ifdefined\listfigurename
  \renewcommand*\listfigurename{List of Figures}
\else
  \newcommand\listfigurename{List of Figures}
\fi
\ifdefined\listtablename
  \renewcommand*\listtablename{List of Tables}
\else
  \newcommand\listtablename{List of Tables}
\fi
\ifdefined\figurename
  \renewcommand*\figurename{Figure}
\else
  \newcommand\figurename{Figure}
\fi
\ifdefined\tablename
  \renewcommand*\tablename{Table}
\else
  \newcommand\tablename{Table}
\fi
}
\@ifpackageloaded{float}{}{\usepackage{float}}
\floatstyle{ruled}
\@ifundefined{c@chapter}{\newfloat{codelisting}{h}{lop}}{\newfloat{codelisting}{h}{lop}[chapter]}
\floatname{codelisting}{Listing}
\newcommand*\listoflistings{\listof{codelisting}{List of Listings}}
\makeatother
\makeatletter
\makeatother
\makeatletter
\@ifpackageloaded{caption}{}{\usepackage{caption}}
\@ifpackageloaded{subcaption}{}{\usepackage{subcaption}}
\makeatother
\makeatletter
\@ifpackageloaded{tcolorbox}{}{\usepackage[skins,breakable]{tcolorbox}}
\makeatother
\makeatletter
\@ifundefined{shadecolor}{\definecolor{shadecolor}{rgb}{.97, .97, .97}}{}
\makeatother
\makeatletter
\makeatother
\makeatletter
\ifdefined\Shaded\renewenvironment{Shaded}{\begin{tcolorbox}[boxrule=0pt, breakable, enhanced, frame hidden, interior hidden, sharp corners]}{\end{tcolorbox}}\fi
\makeatother
\usepackage{bookmark}
\IfFileExists{xurl.sty}{\usepackage{xurl}}{} % add URL line breaks if available
\urlstyle{same}
\hypersetup{
  pdftitle={Running ZZedc Independently for Clinical Research Data Management},
  pdfauthor={Ronald `Ryy' G. Thomas},
  colorlinks=true,
  linkcolor={blue},
  filecolor={Maroon},
  citecolor={Blue},
  urlcolor={Blue},
  pdfcreator={LaTeX via pandoc}}


\title{Running ZZedc Independently for Clinical Research Data
Management}
\usepackage{etoolbox}
\makeatletter
\providecommand{\subtitle}[1]{% add subtitle to \maketitle
  \apptocmd{\@title}{\par {\large #1 \par}}{}{}
}
\makeatother
\subtitle{Evaluating open-source options for investigator-controlled EDC
deployment}
\author{Ronald `Ryy' G. Thomas}
\date{2025-12-07}
\begin{document}
\maketitle

\renewcommand*\contentsname{Table of contents}
{
\setcounter{tocdepth}{3}
\tableofcontents
}

\begin{figure}[H]

{\centering 
\includegraphics[width=0.8\linewidth,height=\textheight,keepaspectratio]{media/images/hero.jpg}

}

\caption{A secure server room with clinical research data management
infrastructure visible}

\end{figure}%

\emph{Investigator-controlled electronic data capture returns data
ownership to the research team.}

\section{Introduction}\label{introduction}

I did not really know how achievable investigator independence in
clinical data management was until I deployed ZZedc on a personal AWS
instance and ran a pilot study without vendor involvement. For years, my
research programmes depended on commercial EDC systems with annual
licensing fees, proprietary data formats, and vendor-controlled
infrastructure. When a vendor relationship deteriorated, the cost of
migration was staggering.

The experience taught me that the barriers to investigator-managed EDC
are not primarily technical. ZZedc installs in minutes, runs on standard
infrastructure, and stores data in portable SQLite databases. The
barriers are organisational: convincing yourself and your institution
that you can manage the infrastructure, security, and maintenance that
vendors handle on your behalf.

This post examines ZZedc as an open-source EDC system designed for
investigator-controlled deployment. It presents three deployment paths
(solo researcher, team research, enterprise), evaluates the technical
architecture, and provides step-by-step guides for each scenario. The
assessment is honest about both the advantages and the significant
responsibilities that come with managing your own clinical data
infrastructure.

\subsection{Motivations}\label{motivations}

\begin{itemize}
\tightlist
\item
  My research programme was paying \$40,000 per year in EDC licensing
  for a system that restricted customisation and locked data in a
  proprietary format.
\item
  A vendor relationship deteriorated, and the estimated cost of data
  migration was \$20,000 -- for data I had collected.
\item
  I wanted to integrate EDC data directly with my R analysis pipelines,
  but the commercial system's export format required manual
  reformatting.
\item
  Multi-site collaborators needed web-accessible data entry, but each
  site's IT department had different policies about vendor-managed
  systems.
\item
  I was curious whether the R/Shiny ecosystem was mature enough to
  support a production EDC application with proper security and
  compliance features.
\end{itemize}

\subsection{Objectives}\label{objectives}

\begin{enumerate}
\def\labelenumi{\arabic{enumi}.}
\tightlist
\item
  Evaluate ZZedc's technical architecture and deployment options for
  investigator-controlled clinical data management.
\item
  Present three deployment paths (solo researcher, team research,
  enterprise) with realistic cost and staffing estimates.
\item
  Provide step-by-step deployment guides for local laptop and AWS server
  installations.
\item
  Honestly assess the tradeoffs between investigator-managed and
  vendor-managed EDC systems, including the hidden costs of
  self-management.
\end{enumerate}

I am documenting my learning process here. If you spot errors or have
better approaches, please let me know.

\begin{center}
\includegraphics[width=0.3\linewidth,height=\textheight,keepaspectratio]{media/images/aws-logo.png}
\end{center}

\emph{Investigator independence requires accepting responsibility for
infrastructure that vendors normally manage.}

\section{Prerequisites and Setup}\label{prerequisites-and-setup}

To deploy ZZedc, you will need one of the following depending on your
deployment path:

\textbf{Solo researcher (local):}

\begin{itemize}
\tightlist
\item
  R (4.0 or later)
\item
  The \texttt{zzedc} R package
\end{itemize}

\textbf{Team research (AWS):}

\begin{itemize}
\tightlist
\item
  AWS account with CLI configured
\item
  Docker
\item
  A registered domain name (for HTTPS)
\item
  Basic command-line familiarity
\end{itemize}

\textbf{Background assumed:} Familiarity with clinical research data
management concepts (data entry, validation, audit trails). No prior
experience with AWS, Docker, or R/Shiny deployment is required for the
solo researcher path.

\section{What is Investigator-Controlled
EDC?}\label{what-is-investigator-controlled-edc}

Investigator-controlled EDC means that the research team, rather than a
vendor, manages the electronic data capture infrastructure. The team
controls where data is stored, how it is accessed, and what happens to
it at the end of the study. This contrasts with vendor-managed systems
where data resides on vendor infrastructure under vendor terms.

Think of it as the difference between renting and owning. With a
vendor-managed system, you rent access to a data management platform: it
is maintained for you, but you do not control the infrastructure, and
your data is stored according to the vendor's policies. With
investigator-managed EDC, you own the infrastructure: you control every
aspect, but you are responsible for maintenance, security, and upkeep.

The decision between these approaches is not primarily technical. It
depends on your team's technical capacity, your institution's compliance
requirements, and the balance between licensing costs and staff time for
self-management.

\section{Getting Started}\label{getting-started}

\subsection{Comparing EDC Approaches}\label{comparing-edc-approaches}

{\def\LTcaptype{none} % do not increment counter
\begin{longtable}[]{@{}
  >{\raggedright\arraybackslash}p{(\linewidth - 4\tabcolsep) * \real{0.3125}}
  >{\raggedright\arraybackslash}p{(\linewidth - 4\tabcolsep) * \real{0.3438}}
  >{\raggedright\arraybackslash}p{(\linewidth - 4\tabcolsep) * \real{0.3438}}@{}}
\toprule\noalign{}
\begin{minipage}[b]{\linewidth}\raggedright
Approach
\end{minipage} & \begin{minipage}[b]{\linewidth}\raggedright
Strengths
\end{minipage} & \begin{minipage}[b]{\linewidth}\raggedright
Tradeoffs
\end{minipage} \\
\midrule\noalign{}
\endhead
\bottomrule\noalign{}
\endlastfoot
Commercial EDC & Polished UI, vendor support & Licensing costs, vendor
lock-in \\
Spreadsheets & Low cost, familiar & Poor validation, no audit trail \\
Custom R/Shiny & Maximum flexibility & Ongoing development required \\
Open-source EDC & Standard formats, lower cost & Requires technical
setup \\
\end{longtable}
}

\subsection{ZZedc Architecture}\label{zzedc-architecture}

\begin{verbatim}
+-------------------------------------------+
|  Your Infrastructure (AWS, Local, Hybrid) |
|                                            |
|  +--------------------------------------+ |
|  |  Web Browser (Any Location)          | |
|  |  https://trial.example.org           | |
|  +--------------------------------------+ |
|            | HTTPS (Automatic)             |
|  +--------------------------------------+ |
|  |  Caddy Reverse Proxy                 | |
|  |  (Automatic Let's Encrypt)           | |
|  +--------------------------------------+ |
|            | Reverse Proxy                 |
|  +--------------------------------------+ |
|  |  ZZedc R/Shiny Application           | |
|  |  (Auth, Forms, Reporting)            | |
|  +--------------------------------------+ |
|            | Database                      |
|  +--------------------------------------+ |
|  |  SQLite Database (Standard Format)   | |
|  |  Your Data, Your Control             | |
|  +--------------------------------------+ |
|                                            |
+-------------------------------------------+
\end{verbatim}

Key design points:

\begin{itemize}
\tightlist
\item
  All infrastructure is standard: Docker containers, standard web server
  (Caddy), open-source database (SQLite).
\item
  No proprietary components: the system uses only open-source software.
\item
  Transparent processes: source code is available for review and
  modification.
\item
  Portable data: SQLite databases can be accessed with any SQL tool.
\end{itemize}

\subsection{What ZZedc Includes}\label{what-zzedc-includes}

\begin{itemize}
\tightlist
\item
  \textbf{Electronic data capture with real-time validation:} Enter
  study data with immediate field-level validation.
\item
  \textbf{Role-based access control:} Five built-in roles (Admin, PI,
  Coordinator, Data Manager, Monitor) with configurable permissions.
\item
  \textbf{Reporting:} Enrollment reports, quality control summaries, and
  statistical overviews.
\item
  \textbf{Data quality framework:} Automated checks for missing data,
  outliers, and consistency.
\item
  \textbf{Audit trail:} Every action is logged with user, timestamp, and
  change history.
\item
  \textbf{Data export:} CSV, Excel, SPSS, and R formats supported.
\end{itemize}

\begin{figure}[H]

{\centering \pandocbounded{
\includegraphics[keepaspectratio]{media/images/ambiance2.jpg}}

}

\caption{A server architecture diagram with database, application, and
web layers labelled}

\end{figure}%

\emph{Standard components and portable data formats enable investigator
independence.}

\section{Deployment Paths}\label{deployment-paths}

\subsection{Path 1: Solo Researcher (Local
Laptop)}\label{path-1-solo-researcher-local-laptop}

\textbf{Best for:} Individual investigators, small pilot studies,
initial prototyping.

Install on your personal laptop in ten minutes:

\begin{Shaded}
\begin{Highlighting}[]
\FunctionTok{install.packages}\NormalTok{(}\StringTok{"zzedc"}\NormalTok{)}
\FunctionTok{library}\NormalTok{(zzedc)}
\NormalTok{zzedc}\SpecialCharTok{::}\FunctionTok{init}\NormalTok{()}
\FunctionTok{launch\_zzedc}\NormalTok{()}
\end{Highlighting}
\end{Shaded}

Data lives in a local SQLite file (\texttt{data/zzedc.db}). You are the
only user. Back up by copying the file to cloud storage.

\textbf{Direct infrastructure cost:} Free (only your laptop
electricity).

\textbf{Staff time:} Minimal (you run it when needed).

\subsection{Path 2: Team Research (Single AWS
Server)}\label{path-2-team-research-single-aws-server}

\textbf{Best for:} Research teams at a single institution or
collaborative multi-site trials.

Deploy on AWS EC2 in approximately 15 minutes:

\begin{Shaded}
\begin{Highlighting}[]
\ExtensionTok{./aws\_setup.sh} \DataTypeTok{\textbackslash{}}
  \AttributeTok{{-}{-}region}\NormalTok{ us{-}west{-}2 }\DataTypeTok{\textbackslash{}}
  \AttributeTok{{-}{-}study{-}name} \StringTok{"Depression Treatment Trial"} \DataTypeTok{\textbackslash{}}
  \AttributeTok{{-}{-}study{-}id} \StringTok{"DEPR{-}2025{-}001"} \DataTypeTok{\textbackslash{}}
  \AttributeTok{{-}{-}admin{-}password} \StringTok{"SecurePass123!"} \DataTypeTok{\textbackslash{}}
  \AttributeTok{{-}{-}domain}\NormalTok{ trial.example.org }\DataTypeTok{\textbackslash{}}
  \AttributeTok{{-}{-}instance{-}type}\NormalTok{ t3.medium}
\end{Highlighting}
\end{Shaded}

Multiple team members access via HTTPS. Database backed up
automatically.

\textbf{Direct infrastructure cost:} Approximately \$30-50/month for EC2
instance.

\textbf{Staff time:} Basic maintenance (Docker handles updates, Caddy
handles HTTPS renewal).

\subsection{Path 3: Enterprise /
Multi-Site}\label{path-3-enterprise-multi-site}

\textbf{Best for:} Large NIH-funded studies, pharmaceutical trials,
production deployments.

Deploy across multiple AWS availability zones with load balancing, RDS
database, S3 backup, and CloudWatch monitoring.

\textbf{Direct infrastructure cost:} \$200-500/month depending on data
volume.

\textbf{Staff time:} Automated (infrastructure-as-code, CI/CD pipeline).

\section{Deployment Scenarios}\label{deployment-scenarios}

\subsection{Scenario 1: Solo Researcher
Prototype}\label{scenario-1-solo-researcher-prototype}

Dr.~Jane is a clinical psychologist planning a small depression
treatment study. She wants to test intervention feasibility before
seeking NIH funding.

\begin{Shaded}
\begin{Highlighting}[]
\FunctionTok{install.packages}\NormalTok{(}\StringTok{"zzedc"}\NormalTok{)}
\FunctionTok{library}\NormalTok{(zzedc)}

\NormalTok{zzedc}\SpecialCharTok{::}\FunctionTok{init}\NormalTok{()}

\FunctionTok{launch\_zzedc}\NormalTok{()}
\end{Highlighting}
\end{Shaded}

Jane has a secure, validated EDC system running locally. She can create
forms for baseline, weekly, and endpoint visits; enrol patients and
enter data; generate enrolment reports; export data to Excel for
analysis; and back up by copying a single file to Dropbox.

\textbf{Direct infrastructure cost:} \$0.

\textbf{Indirect costs:} Staff time for system administration, backup
management, security updates (estimated 2-4 hours/month for solo study).

\subsection{Scenario 2: Multi-Site Trial
Migration}\label{scenario-2-multi-site-trial-migration}

The ADHD research consortium at 5 universities currently uses an
expensive commercial EDC that costs \$40,000/year and is inflexible.
They want to migrate to something more affordable and customisable.

\begin{Shaded}
\begin{Highlighting}[]
\ExtensionTok{./aws\_setup.sh} \DataTypeTok{\textbackslash{}}
  \AttributeTok{{-}{-}region}\NormalTok{ us{-}west{-}2 }\DataTypeTok{\textbackslash{}}
  \AttributeTok{{-}{-}study{-}name} \StringTok{"Multisite ADHD Trial"} \DataTypeTok{\textbackslash{}}
  \AttributeTok{{-}{-}study{-}id} \StringTok{"ADHD{-}MULTI{-}2025"} \DataTypeTok{\textbackslash{}}
  \AttributeTok{{-}{-}admin{-}password} \StringTok{"SecurePassword123!"} \DataTypeTok{\textbackslash{}}
  \AttributeTok{{-}{-}domain}\NormalTok{ adhd{-}trial.org }\DataTypeTok{\textbackslash{}}
  \AttributeTok{{-}{-}instance{-}type}\NormalTok{ t3.large}
\end{Highlighting}
\end{Shaded}

A single centralised instance accessible from all 5 sites with HTTPS and
automatic security certificates. Role-based access: 2 administrators, 5
principal investigators, 20 coordinators, 5 data managers.

\textbf{Direct infrastructure cost:} Approximately \$40-80/month (AWS
EC2 instance + storage).

\textbf{Indirect costs:} Staff time for deployment (20-40 hours initial
setup), ongoing administration and monitoring (5-10 hours/month).

\subsection{Scenario 3: Individual Investigator
Independence}\label{scenario-3-individual-investigator-independence}

Dr.~Robert was running a study with a Biostatistics Lab that managed his
EDC. The lab relationship deteriorated due to cost escalation and
support issues. He wants to take control.

\begin{enumerate}
\def\labelenumi{\arabic{enumi}.}
\tightlist
\item
  Deploy ZZedc independently on AWS.
\item
  Migrate data from the old system (CSV export, import into ZZedc).
\item
  Complete the study independently.
\item
  Archive data in standard SQLite format to institutional repository.
\end{enumerate}

\textbf{Result:} Dr.~Robert owns his data and system. If he works with
another Biostatistics Lab in the future, he can do so without vendor
lock-in.

\section{Security and Compliance}\label{security-and-compliance}

ZZedc includes several security features relevant to regulatory
compliance. Users remain responsible for validating that implementations
meet their institutional requirements.

\textbf{Data governance features:}

\begin{itemize}
\tightlist
\item
  User access logging with timestamps and action records.
\item
  Role-based access control with granular permission assignment.
\item
  Data export capabilities supporting data subject access requests.
\item
  Password encryption and session management.
\end{itemize}

\textbf{Technical security:}

\begin{itemize}
\tightlist
\item
  HTTPS support with automatic certificate management via Let's Encrypt.
\item
  Database encryption options.
\item
  Session timeout controls to prevent unauthorised access.
\end{itemize}

\textbf{Important note:} Regulatory compliance (GDPR, 21 CFR Part 11)
depends on implementation context, institutional policies, and
validation within your specific regulatory environment. Coordinate with
your institutional compliance and IT teams.

\section{Comparison Table}\label{comparison-table}

{\def\LTcaptype{none} % do not increment counter
\begin{longtable}[]{@{}llll@{}}
\toprule\noalign{}
Feature & ZZedc & REDCap & Medidata \\
\midrule\noalign{}
\endhead
\bottomrule\noalign{}
\endlastfoot
Infrastructure & \textasciitilde\$40-500/mo & \$200-1000/mo &
\$500-2000/mo \\
Staff time & Required & Minimal & Minimal \\
Data format & SQLite & Proprietary & Proprietary \\
Customisation & Full source & Limited & Limited \\
Support & Community & Commercial & Commercial \\
Data control & Local & Vendor & Vendor \\
Learning curve & Moderate & Low & Low \\
\end{longtable}
}

\subsection{Things to Watch Out For}\label{things-to-watch-out-for}

\begin{enumerate}
\def\labelenumi{\arabic{enumi}.}
\tightlist
\item
  \textbf{Infrastructure management is real work.} The licensing savings
  from self-managed EDC are offset by staff time for deployment,
  updates, security patching, and troubleshooting. Budget 5-10 hours per
  month for ongoing administration.
\item
  \textbf{Regulatory compliance is your responsibility.} Commercial
  vendors provide compliance frameworks; with self-managed EDC, you must
  validate compliance with your own institutional and regulatory
  requirements.
\item
  \textbf{SQLite has concurrency limitations.} For studies with many
  simultaneous data entry users, SQLite's single-writer constraint may
  cause performance issues. The enterprise path addresses this with
  PostgreSQL.
\item
  \textbf{Domain and DNS configuration adds complexity.} The AWS
  deployment requires a registered domain and DNS configuration, which
  involves coordination with domain registrars and propagation delays.
\item
  \textbf{Backup verification is essential.} Automated backups are only
  useful if you regularly verify that they can be restored. Schedule
  quarterly restore tests.
\end{enumerate}

\begin{figure}[H]

{\centering \pandocbounded{\includegraphics[keepaspectratio]{media/images/ambiance3.jpg}}

}

\caption{A long-term planning document for clinical research
infrastructure beside a laptop showing deployment monitoring}

\end{figure}%

\emph{Long-term sustainability requires planning for maintenance,
updates, and eventual data archival.}

\subsection{Lessons Learnt}\label{lessons-learnt}

\subsubsection{Conceptual Understanding}\label{conceptual-understanding}

\begin{itemize}
\tightlist
\item
  The barriers to investigator-managed EDC are primarily organisational,
  not technical. ZZedc installs and runs with minimal technical effort;
  the challenge is accepting responsibility for infrastructure.
\item
  Total cost of ownership includes both infrastructure costs and staff
  time. The financial comparison with vendor-managed systems depends on
  local staff costs versus vendor licensing.
\item
  Data portability is a more valuable feature than it appears initially.
  When a vendor relationship ends, portable data formats (SQLite, CSV)
  prevent costly migration projects.
\item
  Open-source EDC requires a different support model: community forums
  and documentation replace vendor helpdesks.
\end{itemize}

\subsubsection{Technical Skills}\label{technical-skills}

\begin{itemize}
\tightlist
\item
  The ZZedc \texttt{init()} function generates a project-specific
  configuration through interactive prompts, requiring no prior R/Shiny
  deployment experience.
\item
  The \texttt{aws\_setup.sh} script automates EC2 instance creation,
  Docker deployment, and Caddy HTTPS configuration.
\item
  SQLite databases can be queried directly from R using the
  \texttt{RSQLite} package, enabling tight integration between data
  capture and analysis.
\item
  Caddy's automatic Let's Encrypt integration eliminates the complexity
  of manual HTTPS certificate management.
\end{itemize}

\subsubsection{Gotchas and Pitfalls}\label{gotchas-and-pitfalls}

\begin{itemize}
\tightlist
\item
  The first deployment takes longer than the documented estimates
  because of DNS propagation delays and AWS account setup.
\item
  SQLite's single-writer limitation is not apparent in testing but
  surfaces under production load with multiple simultaneous users.
\item
  Institutional IT departments may resist investigator-managed
  infrastructure; engage them early in the planning process.
\item
  Security patching is an ongoing responsibility that is easy to
  neglect; schedule monthly maintenance windows.
\end{itemize}

\subsection{Limitations}\label{limitations}

\begin{itemize}
\tightlist
\item
  ZZedc requires technical staff for deployment and maintenance.
  Programmes without IT support capacity should consider vendor-managed
  alternatives.
\item
  The comparison with commercial systems (REDCap, Medidata) is
  incomplete because it does not account for the value of
  vendor-provided training, documentation, and support.
\item
  SQLite is appropriate for small to medium studies but may not scale to
  large multi-site trials with hundreds of concurrent users.
\item
  Regulatory compliance validation must be performed independently for
  each implementation context; ZZedc does not provide pre-validated
  compliance frameworks.
\item
  This post evaluates a single open-source EDC system; alternatives such
  as OpenClinica deserve independent evaluation.
\item
  The author has involvement with ZZedc development, which shapes the
  perspective presented here.
\end{itemize}

\subsection{Opportunities for
Improvement}\label{opportunities-for-improvement}

\begin{enumerate}
\def\labelenumi{\arabic{enumi}.}
\tightlist
\item
  Develop a comprehensive compliance validation framework that maps
  ZZedc features to 21 CFR Part 11 and GDPR requirements.
\item
  Create a migration toolkit that automates data import from REDCap,
  Medidata, and CSV sources.
\item
  Build a monitoring dashboard that alerts administrators to security
  issues, backup failures, and performance degradation.
\item
  Implement PostgreSQL support for the enterprise deployment path to
  address SQLite's concurrency limitations.
\item
  Develop institutional deployment guides that address common IT
  department concerns about investigator-managed infrastructure.
\item
  Create a cost comparison calculator that accounts for both licensing
  and staff time across different EDC approaches.
\end{enumerate}

\section{Wrapping Up}\label{wrapping-up}

Running ZZedc independently for clinical research data management is
technically achievable and financially attractive for programmes with
the technical capacity to manage their own infrastructure. The system
installs quickly, stores data in portable formats, and provides the core
features -- data entry, validation, audit trails, role-based access --
that clinical research requires.

What I learnt most from this evaluation is that the decision between
investigator-managed and vendor-managed EDC is not about technology. It
is about organisational capacity. If you have staff who can manage AWS
instances, apply security patches, and verify backups, the financial
case for self-managed EDC is strong. If you do not, the hidden costs of
self-management may exceed the visible costs of vendor licensing.

I encourage investigators interested in this approach to start with the
solo researcher path. Deploy ZZedc on your laptop, create a small pilot
study, and evaluate whether the system meets your needs before
committing to AWS deployment. The migration path from local to cloud is
straightforward.

Main takeaways:

\begin{itemize}
\tightlist
\item
  ZZedc provides investigator-controlled EDC with standard data formats
  and open-source code.
\item
  Three deployment paths serve different scales: solo researcher, team
  research, and enterprise.
\item
  Total cost of ownership includes both infrastructure costs and staff
  time.
\item
  Start with a local pilot before committing to cloud deployment.
\end{itemize}

\section{See Also}\label{see-also}

\textbf{Related posts:}

\begin{itemize}
\tightlist
\item
  \href{../40-testingfordataanalysisworkflow/}{Testing Data Analysis
  Workflows} -- Testing strategies for research pipelines
\end{itemize}

\textbf{Key resources:}

\begin{itemize}
\tightlist
\item
  \href{https://github.com/rgt47/zzedc}{ZZedc GitHub} -- Source code and
  documentation
\item
  \href{https://github.com/rgt47/zzedc/tree/main/\%20deployment}{ZZedc
  Deployment Guide} -- AWS and Docker deployment instructions
\item
  \href{https://www.project-redcap.org/}{REDCap} -- The most widely used
  academic EDC system
\item
  \href{https://www.fda.gov/regulatory-information/\%20search-fda-guidance-documents}{21
  CFR Part 11 Guidance} -- FDA electronic records regulation
\end{itemize}

\section{Reproducibility}\label{reproducibility}

ZZedc is an R package available on GitHub. To install and evaluate:

\begin{Shaded}
\begin{Highlighting}[]
\FunctionTok{install.packages}\NormalTok{(}\StringTok{"zzedc"}\NormalTok{)}
\FunctionTok{library}\NormalTok{(zzedc)}
\NormalTok{zzedc}\SpecialCharTok{::}\FunctionTok{init}\NormalTok{()}
\FunctionTok{launch\_zzedc}\NormalTok{()}
\end{Highlighting}
\end{Shaded}

For AWS deployment:

\begin{Shaded}
\begin{Highlighting}[]
\FunctionTok{git}\NormalTok{ clone https://github.com/rgt47/zzedc.git}
\BuiltInTok{cd}\NormalTok{ zzedc/deployment}
\FunctionTok{chmod}\NormalTok{ +x aws\_setup.sh}
\ExtensionTok{./aws\_setup.sh} \AttributeTok{{-}{-}help}
\end{Highlighting}
\end{Shaded}

\textbf{Project files:}

\begin{verbatim}
zzedcindependence/
  analysis/report/index.qmd  -- This post
  analysis/media/images/      -- Hero, ambiance
\end{verbatim}

\section{Let's Connect}\label{lets-connect}

\begin{itemize}
\tightlist
\item
  \textbf{GitHub:} \href{https://github.com/rgt47}{rgt47}
\item
  \textbf{Twitter/X:} \href{https://twitter.com/rgt47}{@rgt47}
\item
  \textbf{LinkedIn:} \href{https://linkedin.com/in/rgthomaslab}{Ronald
  Glenn Thomas}
\item
  \textbf{Email:} \href{https://rgtlab.org/contact}{rgtlab.org/contact}
\end{itemize}

I would enjoy hearing from you if:

\begin{itemize}
\tightlist
\item
  You spot an error or a better approach to any of the code in this
  post.
\item
  You have suggestions for topics you would like to see covered.
\item
  You want to discuss R programming, data science, or reproducible
  research.
\item
  You have questions about anything in this tutorial.
\item
  You just want to say hello and connect.
\end{itemize}




\end{document}
